\section{Fixed-design finite-batch confirmation with two-layer
uncertainty}\label{confirmation-fixed-design-finite-batch-confirmation-with-two-layer-uncertainty}

29 October 2026. This run increases precision at the ten previously
selected parameter cells without changing the prediction rule, target
window, batch sizes or step size. It is independent of the exploratory
pilot; the pilot observations are excluded from the intervals below.

\subsection{1. Fixed design}\label{confirmation-fixed-design}

\begin{itemize}
\tightlist
\item
  Actual entropy SGD: w=2, nu=1, eta=.1, x0=0, safe interval
  {[}-.25,.25{]}, and the same common-risk construction.
\item
  Cells: n in \{128,512\}, t in \{-1,-.5,0,.5,1\}, with theta=t/(c
  sqrt(n)).
\item
  Independent calibration budget M=round(c\^{}2 n): 10,702 and 42,807
  respectively. The Gaussian comparison is K=r=1, up to integer budget
  rounding.
\item
  Four frozen rules: constant .5, raw moment, window-clipped moment, and
  the previously certified Gaussian step rule. The selected rule is
  recorded verbatim and hashed in \texttt{frozen\_rule.json} and
  \texttt{design.json}.
\item
  Per cell: 16,384 independent reference trajectories and 4,096
  independent calibration repetitions. Calibration samples are shared
  across rules within a repetition, permitting paired comparisons.
  Reference and calibration random streams are independent.
\item
  The trajectory cap is fixed at 1,500 updates. No extension or
  outcome-dependent stopping of simulation is used.
\item
  The design and error allocation were written before the new samples
  were generated. This is a fixed-design confirmation following an
  exploratory pilot, not an externally preregistered study.
\end{itemize}

All updates use sampled nonlinear entropy gradients. Samples are
generated from the actual continuous tilted density by rejection
sampling. The predictor receives only its calibration scores and common
constants, never the generating t, theta, true root, future branch or
reference sample.

\subsection{2. Simultaneous uncertainty
construction}\label{confirmation-simultaneous-uncertainty-construction}

Let C=10 be the number of cells. Allocate alpha\_ref=.025 to reference
trajectories and alpha\_cal=.025 to calibration repetitions.

\subsubsection{2.1. Reference probability, including possible
censoring}\label{confirmation-reference-probability-including-possible-censoring}

At the fixed cap, let p\_R be the probability of a right exit by the cap
and p\_NL the probability of not having exited left by the cap. The
eventual right-exit probability satisfies

\[
p_R\le\pi_n\le p_{NL}.
\]

For each cell use a one-sided exact Clopper--Pearson lower bound on p\_R
and upper bound on p\_NL, each with error alpha\_ref/(2C). Their
endpoints {[}l\_i,u\_i{]} jointly cover every cell's eventual pi\_n with
probability at least 1-alpha\_ref. This remains valid if a cell has
censored paths; lack of an observed censor is not treated as proof that
its population censoring probability is zero.

Set phat\_i=N\_R/N and e\_i=max(phat\_i-l\_i,u\_i-phat\_i). On the
reference event, \textbar phat\_i-pi\_i\textbar\textless=e\_i.

\subsubsection{2.2. Calibration loss uncertainty conditional on the
reference
sample}\label{confirmation-calibration-loss-uncertainty-conditional-on-the-reference-sample}

For each rule and cell, score calibration repetition j by
L\_j=\textbar forecast\_j-phat\_i\textbar. Conditional on the
independent reference sample, these are iid in {[}0,1{]}. Paired
baseline-minus-optimized loss differences are iid in {[}-1,1{]}.

For m observations taking values in an interval of width B, with
unbiased sample variance S\_m\^{}2, use the two-sided empirical
Bernstein radius

\[
b_m(B,\delta)=\sqrt{\frac{2S_m^2\log(4/\delta)}m}
+\frac{7B\log(4/\delta)}{3(m-1)}.
\]

This follows by scaling and applying both tails of Maurer--Pontil
(2009), Theorem 4. Allocate delta=alpha\_cal/(7C) to each of 40
rule-risk estimates and 30 paired gains. Their conditional coverage is
at least 1-alpha\_cal by a union bound; integration over the reference
sample preserves that bound. Independence between the different rules is
not required.

For each rule risk, add e\_i to its calibration radius. For a paired
gain add 2e\_i, by the reverse triangle inequality. For the
deterministic constant rule, propagate {[}l\_i,u\_i{]} through
\textbar.5-p\textbar{} exactly instead of using an unnecessary
calibration bound.

The resulting risk and gain intervals have \textbf{simultaneous coverage
at least 95\%} over all reported cells and rules, under the declared iid
simulation experiment. They include both uncertainty layers. They are
not just conditional standard errors or pointwise normal intervals.
Evaluation of beta quantiles uses SciPy floating-point routines; these
are statistical confidence statements, distinct from the previous
directed-arithmetic Gaussian decision certificates.

\subsubsection{2.3. Finite-grid maximum
risks}\label{confirmation-finite-grid-maximum-risks}

For each n let R\_rule\^{}grid be the maximum of the five true expected
calibration losses at the fixed t cells. If cellwise intervals are
{[}L\_i,U\_i{]}, then

\[
\max_iL_i\le R_{\rm rule}^{\rm grid}\le\max_iU_i.
\]

For the improvement in grid maximum risk, subtract intervals: its lower
bound is max\_i L\_baseline,i minus max\_i U\_optimized,i. The same
simultaneous event covers all these derived comparisons. This is a
five-cell statement, not a continuous-window maximum or finite-n minimax
theorem.

Reference: Andreas Maurer and Massimiliano Pontil, \emph{Empirical
Bernstein Bounds and Sample Variance Penalization}, COLT 2009, Theorem
4. https://www.cs.mcgill.ca/\textasciitilde colt2009/papers/012.pdf

\subsection{3. Results}\label{confirmation-results}

The completed run's tables are appended below. No data-dependent
refinement of the Gaussian rule or cell set is performed.

\subsubsection{3.1. Grid maximum MAE}\label{confirmation-grid-maximum-mae}

All 163,840 reference trajectories exited within the fixed cap; none
were censored. There were 40,960 independent calibration repetitions in
total, each evaluated by all four rules. The intervals below are derived
from one simultaneous coverage event of at least 95\%, including both
sampling layers, and are rounded outward.

{\def\LTcaptype{table} % do not increment counter
\begin{longtable}[]{@{}rlrr@{}}
\toprule\noalign{}
n & Rule & Estimated grid maximum MAE & Simultaneous interval \\
\midrule\noalign{}
\endhead
\bottomrule\noalign{}
\endlastfoot
128 & constant & 0.3426 & {[}0.3338, 0.3511{]} \\
128 & moment & 0.2481 & {[}0.2212, 0.2750{]} \\
128 & clipped & 0.2238 & {[}0.1989, 0.2487{]} \\
128 & optimized & 0.1663 & {[}0.1412, 0.1914{]} \\
512 & constant & 0.3471 & {[}0.3384, 0.3556{]} \\
512 & moment & 0.2476 & {[}0.2206, 0.2746{]} \\
512 & clipped & 0.2228 & {[}0.1979, 0.2478{]} \\
512 & optimized & 0.1684 & {[}0.1434, 0.1933{]} \\
\end{longtable}
}

\subsubsection{3.2. Improvement in grid maximum
risk}\label{confirmation-improvement-in-grid-maximum-risk}

Positive values favor the optimized rule. These compare maxima of
expected losses on the five fixed cells; they do not assert pointwise
dominance.

{\def\LTcaptype{table} % do not increment counter
\begin{longtable}[]{@{}rlrr@{}}
\toprule\noalign{}
n & Baseline & Estimated reduction & Simultaneous interval \\
\midrule\noalign{}
\endhead
\bottomrule\noalign{}
\endlastfoot
128 & constant & 0.1763 & {[}0.1425, 0.2099{]} \\
128 & moment & 0.0819 & {[}0.0299, 0.1338{]} \\
128 & clipped & 0.0576 & {[}0.0076, 0.1075{]} \\
512 & constant & 0.1787 & {[}0.1451, 0.2121{]} \\
512 & moment & 0.0792 & {[}0.0273, 0.1311{]} \\
512 & clipped & 0.0545 & {[}0.0046, 0.1043{]} \\
\end{longtable}
}

Every displayed reduction has a positive lower endpoint, including the
comparisons with window clipping. The result remains model-specific and
does not validate finite-n optimality over all predictors.

\subsubsection{3.3. Cellwise results}\label{confirmation-cellwise-results}

{\def\LTcaptype{table} % do not increment counter
\begin{longtable}[]{@{}
  >{\raggedleft\arraybackslash}p{(\linewidth - 12\tabcolsep) * \real{0.1429}}
  >{\raggedleft\arraybackslash}p{(\linewidth - 12\tabcolsep) * \real{0.1429}}
  >{\raggedleft\arraybackslash}p{(\linewidth - 12\tabcolsep) * \real{0.1429}}
  >{\raggedleft\arraybackslash}p{(\linewidth - 12\tabcolsep) * \real{0.1429}}
  >{\raggedleft\arraybackslash}p{(\linewidth - 12\tabcolsep) * \real{0.1429}}
  >{\raggedleft\arraybackslash}p{(\linewidth - 12\tabcolsep) * \real{0.1429}}
  >{\raggedleft\arraybackslash}p{(\linewidth - 12\tabcolsep) * \real{0.1429}}@{}}
\toprule\noalign{}
\begin{minipage}[b]{\linewidth}\raggedleft
n
\end{minipage} & \begin{minipage}[b]{\linewidth}\raggedleft
t
\end{minipage} & \begin{minipage}[b]{\linewidth}\raggedleft
Reference right probability
\end{minipage} & \begin{minipage}[b]{\linewidth}\raggedleft
Simultaneous reference interval
\end{minipage} & \begin{minipage}[b]{\linewidth}\raggedleft
Moment MAE
\end{minipage} & \begin{minipage}[b]{\linewidth}\raggedleft
Clipped MAE
\end{minipage} & \begin{minipage}[b]{\linewidth}\raggedleft
Optimized MAE
\end{minipage} \\
\midrule\noalign{}
\endhead
\bottomrule\noalign{}
\endlastfoot
128 & -1 & 0.16321 & {[}0.1545, 0.1721{]} & 0.18038 & 0.12645 &
0.15757 \\
128 & -0.5 & 0.30817 & {[}0.2973, 0.3192{]} & 0.23400 & 0.20093 &
0.16552 \\
128 & 0 & 0.49817 & {[}0.4863, 0.5101{]} & 0.24813 & 0.22381 &
0.16306 \\
128 & 0.5 & 0.69037 & {[}0.6793, 0.7013{]} & 0.23015 & 0.19765 &
0.16299 \\
128 & 1 & 0.84259 & {[}0.8338, 0.8511{]} & 0.18340 & 0.13257 &
0.16626 \\
512 & -1 & 0.16205 & {[}0.1534, 0.1710{]} & 0.18425 & 0.13184 &
0.16259 \\
512 & -0.5 & 0.30695 & {[}0.2960, 0.3180{]} & 0.22602 & 0.19352 &
0.16187 \\
512 & 0 & 0.50385 & {[}0.4920, 0.5157{]} & 0.24760 & 0.22282 &
0.16288 \\
512 & 0.5 & 0.68506 & {[}0.6739, 0.6960{]} & 0.23076 & 0.19884 &
0.16429 \\
512 & 1 & 0.84711 & {[}0.8384, 0.8556{]} & 0.18185 & 0.13448 &
0.16837 \\
\end{longtable}
}

The interval construction supports positive paired improvements over the
raw moment rule in 6/10 cells and over the clipped rule in 2/10 cells.
At endpoint cells clipping is often better, so there is no claim of
uniform pointwise improvement. At t=0 the target equals .5 exactly by
symmetry in this family; for consistency of the primary analysis, that
cell was scored with the independent reference estimate, without using
the exact value to narrow its interval.

The largest observed exit step was 720. Raw exit arrays, forecast
arrays, fixed-design metadata and the frozen rule are retained. The
checksum of the rule and the aggregate counts and MAEs were
independently checked against the retained NPZ files.

\subsection{4. Recommended use in the
manuscript}\label{confirmation-recommended-use-in-the-manuscript}

Use the grid-risk reduction table as finite-batch confirmation of a
frozen prospective predictor. Keep the Gaussian minimax certificates in
the theory/numerics section and the two-layer confidence construction in
the experimental appendix. The new experiment does not test unknown
nuisance directions; that extension is presently theoretical. Full
visual adaptation, unknown-baseline geometry, and a finite-n
continuous-window guarantee remain outside this evidence.
